\section{Bölüm sonu çözümlü problemler}

\begin{enumerate}
  \item \textbf{Merkez ve yayılma.} Beş öğrencinin oturum sayıları $2,3,3,4,8$'dir. Ortalama, medyan, ranj ve örneklem standart sapmasını hesaplayın.

  \textbf{Çözüm:} Ortalama $\bar{x}=20/5=4$, medyan 3 ve ranj $8-2=6$'dır. Kareli sapmalar toplamı $(-2)^2+(-1)^2+(-1)^2+0^2+4^2=22$ olduğundan örneklem varyansı $22/(5-1)=5{,}5$, standart sapma $s=\sqrt{5{,}5}\approx2{,}35$ olur.

  \item \textbf{Çeyrekler ve aykırı değer.} Sıralı puanlar $4,5,6,6,7,8,9,15$ olsun. Ortanca-yarılar yöntemiyle $Q_1$, $Q_3$, IQR ve üst aykırı değer sınırını bulun.

  \textbf{Çözüm:} Alt yarının medyanı $Q_1=(5+6)/2=5{,}5$, üst yarının medyanı $Q_3=(8+9)/2=8{,}5$'tir. $IQR=3$ ve üst sınır $8{,}5+1{,}5(3)=13$ olur. 15 aykırı değer adayıdır; otomatik olarak silinmez.

  \item \textbf{Sınıflandırılmış ortalama.} Puan sınıflarının orta noktaları 15, 25 ve 35; frekansları sırasıyla 4, 10 ve 6'dır. Yaklaşık ortalamayı bulun.

  \textbf{Çözüm:} Toplam frekans $n=20$ ve ağırlıklı toplam $4(15)+10(25)+6(35)=520$'dir. Sınıflandırılmış ortalama $\bar{x}_g=520/20=26$ olur. Bu değer, her sınıftaki gözlemlerin orta noktada bulunduğu varsayımına dayanan yaklaşık bir özettir.
\end{enumerate}